\section{Literature Review}
\label{sec:litreview}

\subsection{Classical Portfolio Optimization}
\label{subsec:lit_classical}

The foundation of quantitative portfolio management is the mean-variance framework
introduced by \citet{Markowitz1952}, which characterizes optimal portfolios as those
lying on the efficient frontier---the set of allocations that maximize expected return
for a given level of variance, or equivalently minimize variance for a target return.
Despite its theoretical elegance, mean-variance optimization is notoriously sensitive
to estimation error in the sample mean and covariance matrix, which causes the optimal
weights to be unstable and the out-of-sample performance to deteriorate significantly
\citep{Merton1980}.

\citet{DeMiguel2007} conduct an extensive empirical study of fourteen mean-variance
extensions across seven datasets and find that none consistently outperforms the naive
$1/N$ equal-weight portfolio in terms of Sharpe ratio or turnover out-of-sample.
This result---often called the \emph{1/N puzzle}---suggests that estimation error in
classical methods is severe enough to erode all theoretical benefits from
optimization.
The Minimum-Variance Portfolio (MVP), which ignores the mean and relies solely on the
covariance structure, partially mitigates this problem and tends to outperform
full mean-variance models, though it still underperforms simple equal-weight strategies
in many settings \citep{Clarke2011}.

Risk Parity is a more recent alternative that allocates weights inversely proportional
to individual asset volatility contributions, so that each asset contributes equally to
total portfolio risk \citep{Qian2011, Maillard2010}.
Risk Parity gained practical traction after the 2008 financial crisis because its
diversification principle is independent of return estimates.
However, like MVP and equal-weight, Risk Parity is entirely regime-agnostic: it
applies the same weighting logic regardless of whether the economy is tightening,
easing, or at a neutral stance.
In emerging markets with pronounced monetary cycles, this insensitivity to the
macro environment is a material limitation.

\subsection{Reinforcement Learning in Finance}
\label{subsec:lit_rl_finance}

The application of reinforcement learning to financial decision-making has grown
rapidly since the introduction of deep Q-networks \citep{Mnih2015}, which demonstrated
that neural network function approximators could learn superhuman policies directly
from raw inputs.
Early applications to portfolio management used Q-learning on discretized action
spaces \citep{Moody1998a}, but the continuous nature of portfolio weights motivated a
shift to policy-gradient and actor-critic methods.

\citet{Jiang2017} propose a Convolutional Neural Network-based portfolio management
system trained with a policy-gradient algorithm, demonstrating consistent
outperformance over buy-and-hold and equal-weight strategies on cryptocurrency markets.
\citet{Ye2020} apply the Proximal Policy Optimization (PPO) algorithm to U.S. equity
portfolios and show that including technical indicators in the state space improves
risk-adjusted returns.
Deep Deterministic Policy Gradient (DDPG) \citep{Lillicrap2015} has been widely
adopted for its ability to handle continuous action spaces and has been applied to
stock trading \citep{Xiong2018}, commodity futures \citep{Liang2018}, and
multi-asset allocation \citep{Soleymani2021}.
Soft Actor-Critic (SAC) \citep{Haarnoja2018}, which maximizes a combination of
expected return and entropy, has shown robustness advantages in financial settings
due to its inherent exploration mechanism.

A common thread across these studies is that they operate in an \emph{online} setting,
where the agent simultaneously interacts with the market environment and updates its
policy.
This conflation of training and evaluation creates a risk of look-ahead bias,
particularly in simulation-based studies where the environment is replayed from
historical data and the agent can, in principle, extract information about future
returns from the reward signal during training.

\subsection{Offline Reinforcement Learning and Conservative Q-Learning}
\label{subsec:lit_offline_rl}

Offline RL---also known as batch RL---addresses the look-ahead bias problem by
restricting training to a fixed dataset collected by a \emph{behavioral policy}
that is separate from the policy being learned \citep{Levine2020}.
The agent never interacts with the live environment during training, eliminating
any possibility of exploiting future information.
This paradigm is particularly attractive for finance, where live experimentation is
costly and historical data is abundant.

The central challenge of offline RL is the \emph{distributional shift} problem:
Q-values for actions that are underrepresented in the offline dataset can be
severely overestimated, leading to policies that select actions the behavioral
policy would never take and that may perform catastrophically in deployment.
\citet{Kumar2020} propose Conservative Q-Learning (CQL) to address this.
CQL augments the standard Bellman loss with a regularization term that penalizes
the expected Q-value of actions sampled from the current policy while rewarding
the expected Q-value of actions sampled from the behavioral data.
Formally, this conservatism penalty minimizes a lower bound on the true Q-value
for out-of-distribution actions, preventing the agent from assigning unrealistically
high values to unexplored regions of the action space.
\citet{Kumar2020} demonstrate that CQL achieves state-of-the-art performance on
the D4RL offline RL benchmark \citep{Fu2020}, outperforming prior methods including
Behavior-Cloning, BCQ \citep{Fujimoto2019}, and BEAR \citep{Kumar2019}.

In finance, offline RL has attracted growing attention.
\citet{Yuan2025} apply an offline policy within a bilevel optimization framework
for sequential portfolio optimization, highlighting the practical importance of
ensuring that the offline policy does not merely memorize the optimal actions
from the training data.
The present paper is, to our knowledge, among the first to apply CQL specifically
to multi-asset equity portfolio optimization in the Brazilian market, combining
offline safety guarantees with explicit tail-risk control.

\subsection{CVaR in Portfolio Optimization}
\label{subsec:lit_cvar}

Conditional Value-at-Risk (CVaR), also known as Expected Shortfall or Tail
Value-at-Risk, was formally introduced as a coherent risk measure by
\citet{Rockafellar2000}, who also proposed a linear-programming formulation
that makes CVaR minimization tractable in the context of portfolio optimization.
Unlike Value-at-Risk (VaR), which only characterizes the loss threshold at a
given confidence level but ignores the severity of losses beyond that threshold,
CVaR measures the expected loss conditional on exceeding the VaR threshold.
CVaR satisfies the axioms of coherent risk measures---translation invariance,
monotonicity, positive homogeneity, and subadditivity---making it a theoretically
sound basis for portfolio risk control \citep{Artzner1999}.

The application of CVaR to reinforcement learning reward design has received
increasing attention in the risk-sensitive RL literature.
\citet{Tamar2015} develop a policy-gradient algorithm for CVaR-constrained MDPs,
proving convergence to a local optimum.
More recently, \citet{Dabney2018} propose Distributional RL methods that learn
the full return distribution rather than only its expectation, enabling
quantile-based risk measures to be incorporated directly into the value function.
In the portfolio context, a CVaR reward penalizes the agent for generating
return distributions with heavy left tails, which is particularly relevant in
emerging markets where political risk, currency crises, and commodity shocks
produce return distributions with excess kurtosis and negative skewness.

\subsection{Macroeconomic Regimes and Emerging Markets}
\label{subsec:lit_macro_regimes}

The literature on regime-switching models in finance dates to \citet{Hamilton1989},
who demonstrated that U.S. GDP growth follows a two-state Markov switching process.
Subsequent work extended regime-switching to equity returns, showing that volatility,
correlation structure, and expected returns vary substantially across expansionary and
recessionary regimes \citep{Ang2002}.
In portfolio management, conditioning on regimes has been shown to improve
risk-adjusted performance: \citet{Guidolin2007} find that mean-variance investors
who condition on business-cycle regimes achieve economically large utility gains
relative to unconditional allocations.

The Brazilian market presents a particularly compelling case for regime-aware
strategies.
The SELIC rate---the primary instrument of Brazilian monetary policy---has fluctuated
between 2\% and over 14\% during our sample period (2015--2024), with each major
cycle coinciding with distinct equity-market dynamics.
\citet{Sachsida2006} document that Brazilian asset prices are more responsive to
domestic monetary shocks than to international ones, underscoring the importance of
incorporating SELIC dynamics into portfolio models.
Additionally, the Brazilian market is characterized by relatively high equity
premiums, elevated credit spreads, and persistent currency volatility that amplify
the macro-financial transmission mechanism compared with developed markets
\citep{Bonomo2011}.
These structural features motivate conditioning the allocation on the monetary
stance rather than on price history alone.
They also set a ceiling on what the state vector used here can capture: sovereign
credit risk and implied volatility are distinct dimensions of emerging-market risk,
and neither is observed in this study (Section~\ref{subsec:data_macro}).
